\documentclass[11pt]{article}

%====================
% Basic setup
%====================
\usepackage[left=0.7in, right=0.7in, top=0.7in, bottom=0.7in]{geometry}
\usepackage{hyperref}
\usepackage{xcolor}
\usepackage{enumitem}
\usepackage{titlesec}
\usepackage{setspace}

\hypersetup{
    colorlinks=true,
    urlcolor=black
}

\setlength{\parindent}{0pt}
\setlength{\parskip}{4pt}
\pagenumbering{gobble}






%====================
% Section formatting
%====================
\titleformat{\section}
  {\large\bfseries}
  {}
  {0pt}
  {}
  [\vspace{2pt}\titlerule]

%====================
% Document
%====================
\begin{document}





%====================
% Header
%====================
{\Huge \textbf{Angel Fernando Bórquez Guerrero}} \\
\vspace{-0.4cm}
\par \href{https://www.linkedin.com/in/fernandoborquez/}{LinkedIn} \quad | \quad
\href{https://github.com/FernyBoy}{GitHub} \quad | \quad
\href{mailto:fernandoborquez215@icloud.com}{fernandoborquez215@icloud.com} \quad | \quad
+52 633 100 5991 \\
Hermosillo, Sonora.





%====================
% Education
%====================
\section*{Education}
\textbf{Universidad de Sonora} \hfill CGPA: 91.91/100\\
Bachelor of Science in Computer Science \hfill \textbf{August 2023 -- Present}



%====================
% Professional Experience
%====================
\section*{Professional Experience}
\textbf{TE Connectivity} \hfill Supervisor: M.Sc. Martín Vega \\
AI Engineer Intern \hfill \textbf{June 2026 -- September 2026}
\vspace{-0.3cm}
\begin{itemize}[leftmargin=*]
    \item Designed and developed an AI agent for industrial alert analysis, implementing its orchestration, state management, tool integration, and execution workflows to support automated analysis and decision-making.

    \vspace{-0.3cm}
    \item Re-architected the agent's orchestration using LangGraph, replacing the previous workflow implementation with a graph-based architecture that improved state management, modularity, and execution efficiency.

    \vspace{-0.3cm}
    \item Optimized the agent's architecture and execution pipeline, reducing response times from up to three minutes to near-instant responses and substantially improving its practical usability.

    \vspace{-0.3cm}
    \item Designed and deployed a shared Linux development server for interns, configuring the system and user access to provide a centralized environment for developing and hosting research and engineering projects.
\end{itemize}


\newpage
\vspace{0.2cm}
\textbf{ACARUS — Supercomputing Area} \hfill Supervisor: Dr. María del Carmen Heras Sánchez\\
Professional Intern \hfill \textbf{September 2025 -- Present}
\vspace{-0.3cm}
\begin{itemize}[leftmargin=*]
    \item Professional internship in the Supercomputing area at the University of Sonora.
    
    \vspace{-0.3cm}
    \item Part of a development team building a specialized forum for users of \textit{Yuca}, the university HPC system.
    
    \vspace{-0.3cm}
    \item Implementing features such as user posts, direct messaging, and service request workflows.
    
    \vspace{-0.3cm}
    \item Supporting an active HPC user community by improving communication and technical support processes.
\end{itemize}


%====================
% Research Experience
%====================
\section*{Research Experience}
\textbf{University of Guadalajara} \hfill Guadalajara, México \\
\textit{Undergraduate Research Intern} (Advisor: Dr. Rafael Morales Gamboa) \hfill \textbf{June 2025 -- July 2025}
\vspace{-0.3cm}
\begin{itemize}
    \item Validated and generalized an Entropic Associative Memory (EAM) model across high-dimensional benchmark datasets, culminating in a peer-reviewed publication at \textit{IEEE ENC 2025}.

    \vspace{-0.3cm}
    \item Designed an experimental novelty detection framework to rigorously evaluate the EAM architecture's out-of-distribution (OOD) class rejection performance under varying operational bounds.

    \vspace{-0.3cm}
    \item Coupled deep perceptual networks (convolutional autoencoders and classifiers) to project raw image spaces into bounded latent vector representations compatible with associative memory topologies.

    \vspace{-0.3cm}
    \item Formulated quantitative control protocols using specialized ``no-response'' metrics to model partial memory saturation and characterize systematic false-positive rejection dynamics.

    \vspace{-0.3cm}
    \item Engineered an end-to-end reproducible simulation framework, modularizing class density distributions, balance ratios, and multi-class retrieval evaluation regimes.
\end{itemize}


%====================
% Publications
%====================
\section*{Publications}
\begin{itemize}
    \item \textbf{A. F. Bórquez Guerrero}, R. Morales Gamboa, et al. ``\textit{Large-Scale Validation and Analysis of the Rejection Capacity in Entropic Associative Memory}.'' In \textbf{Proceedings of the 2025 Mexican International Conference on Computer Science (ENC) \& IEEE ICEV}, Orizaba, Veracruz, México, 2025. IEEE. \\
    \footnotesize\textit{Indexing: Online ISSN: 2332-5712 $\vert$ ISBN: 979-8-3315-9026-0/25 $\vert$ IEEE Xplore.}
\end{itemize}


%====================
% Competitions & Projects
%====================
\section*{Competitions \& Projects}
\textbf{TE AI CUP 2026 -- International 1st Place} \hfill Project: Industrial Alert Agent\\
Project Leader \hfill \textbf{January 2026 - June 2026}
\vspace{-0.3cm}
\begin{itemize}[leftmargin=*]
    \item Led a 6-member interdisciplinary team in an international AI competition, guiding the technical development of an AI-driven system for industrial alert analysis that achieved 1st place. 

    \vspace{-0.3cm} 
    \item Developed an AI agent integrating Large Language Models (LLMs) for automated alert analysis and decision support.

    \vspace{-0.3cm} 
    \item Developed an end-to-end data and machine learning pipeline involving data cleaning, text preprocessing, embedding generation, clustering, data validation, and predictive modeling. 

    \vspace{-0.3cm}
    \item Developed a predictive modeling system to estimate the likelihood of future industrial alerts, providing a proactive component that complemented the analysis of historical patterns and enabled earlier identification of potential issues.

    \vspace{-0.3cm} 
    \item Architected the project around clean, standardized, and scalable software structures, establishing conventions for modules, tools, configuration, state management, and component integration. 
\end{itemize}



\textbf{HSIL Hackathon 2026 -- 3rd Place} \hfill Project: ShiftGuardian\\
Application Designer \& AI Agent Architect \hfill \textbf{April 2026}
\vspace{-0.3cm}
\begin{itemize}[leftmargin=*]
    \item Designed the application interface and AI agent architecture for a clinical decision support system supporting nurse handoffs and patient prioritization.
    
    \vspace{-0.3cm}
    \item Integrated a RAG-based LLM workflow concept to summarize clinical reports, identify critical tasks, and reduce cognitive load during shift transitions.
\end{itemize}



\textbf{MICAI 2025 MeDA Challenge -- National 1st Place} \hfill \textbf{October 2025}
\vspace{-0.3cm}
\begin{itemize}[leftmargin=*]
    \item Contributed to a 1st Place victory in a national academic competition focused on Medical Domain Adaptation using Self-Supervised Learning (SSL), leading to an invitation to publish in a special issue of the \textit{Health Information Science and Systems} journal.
    
    \vspace{-0.3cm}
    \item Structured and standardized the dataset and training assets to establish a reproducible data pipeline for efficient model development and experimentation.
    
    \vspace{-0.3cm}
    \item Investigated approaches for cross-domain generalization and representation learning across heterogeneous medical imaging datasets.

    \vspace{-0.3cm}
    \item Developed visual and presentation assets used to communicate the project during the competition and dissemination stages.
\end{itemize}



\textbf{TE AI CUP 2025 -- International 3rd Place} \hfill Project: AI Generated Machine Downtime Insights\\
UI Developer \hfill \textbf{January -- May 2025}
\vspace{-0.3cm}
\begin{itemize}[leftmargin=*]
    \item Designed and developed the UI/UX for an AI-based industrial downtime analysis platform, translating complex machine time-series data and unsupervised learning outputs into clear, actionable visual insights.
    
    \vspace{-0.3cm}
    \item Developed an accessible and user-centered interface for non-technical users, simplifying the exploration and interpretation of AI-generated insights without requiring prior knowledge of machine learning.
    
    \vspace{-0.3cm}
    \item Defined user flows, information hierarchy, and interaction patterns to improve navigation and facilitate the transition from high-level downtime trends to detailed analytical results.

    \vspace{-0.3cm} 
    \item Collaborated with the data science and engineering teams to align the interface with the underlying analytical models and ensure that visual representations accurately reflected model outputs.
\end{itemize}



\end{document}
